% 09-use.tex: holders, relying parties, the broker, and the stranger's path.
\section{Using the credential: holders, relying parties and the login that keeps no record}
\label{sec:use}

\motif{Holding a credential is not the same as having to show all of it.}

\analogy{Townspeople carry their own papers. A merchant learns what the merchant needs and no more, from papers the customer chooses to show, and two merchants comparing their ledgers cannot tell that they served the same person.}

\subsection{The wallet is a protocol, not a product}

The reference holder agent, \code{scripts/polaris-wallet.py}, holds a credential as a file, shows
it, verifies it offline through the detached verifier, presents it, proves membership in zero
knowledge against a published epoch, signs a document on the holder's behalf, delegates a bounded
authority to an agent (Section~\ref{sec:agents}) and logs in to a relying party. Every message it
emits is specified so that any client can emit the same. Native mobile and desktop clients and a
browser bridge are product engineering that a reference implementation on notional data does not
attempt; they would follow the protocol, and never lead it. The physical presentation path, a card
that signs a device's challenge, is the emulator and verifier device of Section~\ref{sec:core}.
\sImpl{} for the protocol and the reference script; \sFut{} for native clients.

A presentation is an unsigned wrapper around signed things: the issuer-signed credential, an
optional stapled status assertion so that authorization is decidable with no connectivity, an
optional membership proof, an optional holder signature over the verifier's nonce, the context and
disclosure level presented under, and an opaque presentation code. The verifier decides it offline
and reports the code as present or absent without ever interpreting it, because a duress code is
matched only by the issuer, out of sight; a duress presentation is byte-for-byte the same shape as a
normal one. For transfer by QR or NFC the presentation is compressed and split into frames that each
name the digest of the whole; a receiver takes them in any order and refuses, before parsing
anything, frames naming different digests, a missing index, a conflicting duplicate, or a
reassembled payload whose digest does not match. The decoder is resource-bounded: at most 9,999
frames, 512~KiB compressed and 2~MiB inflated, so a decompression bomb is refused at the limit rather
than decompressed. The bounds were added after an outside review named the gap. \sImpl

\subsection{The relying party}

The reference relying party (\code{scripts/\allowbreak{}polaris-relying-party.py}) decides accept or reject for a
presentation by combining offline authenticity with online status, refusing a revoked, tampered or
foreign-issuer credential and staying blind to duress, and an end-to-end drill runs the whole
wallet-to-relying-party matrix under real signatures every release. A relying party registers with
an issuer as an organisation with a scope that a schema constraint bounds to exactly \code{verify},
\code{authenticate}, or both, so that a fourth kind of access is a constitutional change rather than
a configuration one. Identity never becomes a login product by accident. A
relying party is also an object about which decisions are recorded: registering one, turning off the
zero-knowledge step-up its holders had to satisfy, disabling or re-enabling it, each writes an
append-only event, a change that weakens what the party must satisfy is flagged as such and refused
without a stated reason, and a change made from a database shell is recorded on the same terms as
one made through the tool, because the writer is a trigger and not the tool. The first version of
the rule that decides what counts as a weakening ranked a capability by the length of its name, so
granting a capability the party did not hold read as harmless; it was found by running it against a
loaded database, not by reading it. \sImpl

\subsection{Logging in without a login record}

The auth broker lets a relying party offer \emph{log in with a Polaris credential} through an
authorization code with PKCE. The holder presents the credential at its issuing authority's
instance by possession, and the relying party exchanges the resulting code for an identity token
signed by the issuing agency, which it verifies offline. Four guards keep this from becoming a login
record. The token's subject is the per-relying-party handle of Section~\ref{sec:privacy}, never a
person identifier and not the same value at every relying party. The token
carries no claim beyond the vocabulary: context, disclosure level, the assurance reached, enrolment
status, instants. The authorize route writes nothing, and the token endpoint's only write is the
code's hash into a register that holds no subject and no relying party. And duress is served
identically. \sImpl

A relying party's demands, a required zero-knowledge step-up or a required enrolment status, bind
the authority and do not travel only in the holder-side request: the policy is registered on the
relying party's row and applied first, a request may add a requirement and never remove one, and a
context other than the registered one is refused. The authorization code is encrypted as well as
signed. And the conformance
contract was later found to verify an identity token's signature and never its audience or nonce, so
a verifier that ignored both, and therefore accepted a token minted for one relying party at another,
passed every case; four cases fixed it. \sImpl

\subsection{The stranger's path}

A reader can reach one accepted presentation from software nobody here wrote, on a clean machine, in about ten
minutes, without cloning the repository: install \code{polaris-oid4vp} from the registry, start the
published wallet image at its pinned digest, generate the verifier's keys, give the wallet a
credential bound to a key it generated itself, run the verifier, present. The page that says how was
executed from an empty directory before it was written down, which caught its own defect on the way
(a buffered pipe that left the verifier's log empty), and it ends with the line that matters:
\code{<- 200 authentic, claims ['cnf', 'family\_name', 'given\_name', 'iat', 'iss', 'vct']}. If it
fails for a reader, the page says, that is the most useful bug the project can receive. Nobody
outside the repository has yet reported walking it. \sImpl{} for the path; \sExt{} for a walker
who is not the author.

\subsection{What is deliberately not built}

There is no session at the authority, no consent screen, no discovery document, and no browser
redirect choreography: a relying party integrates from the wire specification and the registry. A
WebAuthn browser bridge is not built; the holder-side key it would rest on now exists. No relying
party outside the repository has integrated against the native protocol. One outside wallet has
presented over OpenID4VP, and no outside relying party has yet run the verifier that accepted it.
\sExt

\begin{layercard}{Layer card: use}
\qa{What it is}{The wallet protocol and reference script, presentations and QR framing, the reference relying party, relying-party registration with a bounded scope and a recorded history, the auth broker, and the stranger's path to an accepted presentation from an outside wallet.}
\qa{Why it exists}{A credential nobody can hold, present or accept is a database row. This layer is where the credential meets software the issuer does not run.}
\qa{What it trusts}{The signed objects inside a presentation; the issuer's key for the identity token; the relying party's registered policy.}
\qa{What it refuses}{To interpret the presentation code; to write a login record; to accept a scope outside the CHECK; to let a holder's request weaken the relying party's registered demand; to lower a relying party's bar without a reason on record; to decompress an unbounded payload.}
\qa{What it can see}{A relying party sees the credential value and the verdicts, or on the zero-knowledge path a proof and a scoped handle. The authority sees a possession proof and consumes a code hash.}
\qa{What it cannot see}{The authority cannot see who logged in where; the relying party cannot see the person's record.}
\qa{If it fails}{A replayed code hits the consumed register; a wrong presentation, an unmet step-up, a verify-only bearer at the broker, a stale or tampered frame are each refused in the drills; a weakened relying-party bar with no reason is refused by the database.}
\qa{Checked by}{\code{check\_holder\_wallet}, \code{check\_wallet\_presentation}, \code{check\_holder\_verifier\_flow}, \code{check\_auth\_broker}, \code{check\_broker\_policy\_bound}, \code{check\_relying\_party\_api}, \code{check\_qr\_resource\_bounds}, \code{check\_documented\_verifier\_commands\_run}.}
\qa{Now or later}{Now, as a protocol, and now for one outside wallet on one path. Native clients and an outside relying party are later.}
\qa{Inside and outside}{Inside: the privacy boundary. Outside: federation, because a relying party rarely trusts only one authority.}
\end{layercard}

\nextlayer{A relying party that trusts one authority is easy. A society has many authorities, none
of which should be able to speak for the others, and none of which should hold every person's
record. The next layer is how independent authorities recognise one another without collapsing into
one.}
